\scanpage{46}
\begin{center}-- 18 --\end{center}
\begin{flushright}n\textsuperscript{o} 370\end{flushright}

\uline{positive sur} $\mathbf Q_p$ \uline{telle que toute classe suivant} $p^n\mathbf Z_p$ ($n$ \uline{entier} $\geq0$) \uline{soit de mesure} $p^{-n}$.

La restriction de $\alpha$ à $\mathbf Z_p$ est évidemment une mesure de Haar sur $\mathbf Z_p$.

6. \uline{Définition locale d'une mesure de Haar.}

PROPOSITION 7.-- \uline{Soient} $G$ \uline{un groupe localement compact}, $V$ \uline{une partie ouverte de} $G$, $\mu$ \uline{une mesure positive non nulle sur} $V$, \uline{ayant la propriété suivante} : \uline{si} $U$ \uline{est une partie ouverte de} $V$ \uline{et si} $s\in G$ \uline{est tel que} $sU\subset V$, \uline{l'image de la mesure} $\mu_U$, \uline{induite par} $\mu$ \uline{sur} $U$, \uline{par l'homéomorphisme} $x\to sx$ \uline{de} $U$ \uline{sur} $sU$, \uline{est} $\mu_{sU}$. \uline{Alors il existe sur} $G$ \uline{une mesure de Haar à gauche} $\alpha$ \uline{et une seule induisant} $\mu$ \uline{sur} $V$.

Pour tout $s\in G$, soit $\mu_s$ l'image de $\mu$ par l'homéomorphisme $x\to sx$ de $V$ sur $sV$. La restriction de $\mu_s$ à $V\cap sV$ est l'image de $\mu_{s^{-1}(V)\cap V}$ par la restriction de $x\to sx$ à $s^{-1}(V)\cap V$ ; par hypothèse, cette image est $\mu_{V\cap sV}$. Par translation, on en conclut que $\mu_s$ et $\mu_t$ ont même restriction à $sV\cap tV$ quels que soient $s$, $t$. D'après la prop.1 du chap. III, \S\ 3, n\textsuperscript{o} 1, il existe donc une mesure $\alpha$ sur $G$ induisant $\mu_s$ sur $sV$ quel que soit $s$. Il est clair que $\alpha$ est l'unique mesure de Haar à gauche sur $G$ induisant $\mu$ sur $V$.

COROLLAIRE.-- \uline{Soient} $G$, $G'$ \uline{deux groupes localement compacts}, $V$ (resp. $V'$) \uline{un voisinage ouvert de l'élément neutre de} $G$ (resp. $G'$), $\varphi$ \uline{un isomorphisme local de} $G'$ \uline{à} $G$, \uline{défini sur} $V'$, \uline{tel que} $\varphi(V')=V$. \uline{Soient} $\alpha'$ \uline{une mesure de Haar à gauche sur} $G'$, \uline{et} $\alpha'_{V'}$ \uline{sa restriction à} $V'$. \uline{Alors} $\varphi(\alpha'_{V'})$ \uline{est la restriction à} $V$ \uline{d'une mesure de Haar à gauche unique} $\alpha$ \uline{sur} $G$.

Soit $V_1$ un voisinage ouvert de $e$ dans $G$ tel que $V_1V_1^{-1}\subset V$. Soit $\mu$ la restriction de $\varphi(\alpha'_{V'})$ à $V_1$. Soient $U$ une partie ouverte de $V_1$ et $s\in G$ tels que $sU\subset V_1$. On a $s\in V_1V_1^{-1}\subset V$, donc $s=\varphi(s')$ avec un $s'\in V'$. Soit $x\in U$. On a $x=\varphi(x')$ avec un $x'\in V'$, donc $sx=\varphi(s')\varphi(x')=\varphi(s'x')$ puisque $sx\in sU\subset V$. Comme les translations à gauche dans $G'$ conservent $\alpha'$, on voit que $V_1$ et $\mu$ satisfont aux conditions de la prop. 7. Soit $\alpha$ la mesure de Haar à gauche sur $G$ induisant $\mu$ sur
\mspage{45}
